\chapter{Methodology}
\label{ch:methodology}

\section{Analysis principle}

The design separates the original MIT experiment, an exploratory linear follow-up and an
independent-corpus prediction study. The original model configurations and contrasts were
fixed before the perceived-attribute audit. Its findings subsequently motivated the linear
follow-up, which therefore cannot be interpreted as a prospectively independent confirmation.
Audit attributes remained outside ordinary fitting and selection in both analyses. The first
sections describe the original design; the later sections state the changes used in the
follow-up and independent corpus. Their distinct units, outcome scales and summaries are
retained in the Results.

\section{Dataset contract}

\subsection{Units, identifiers and target}

The authorised MIT Interview distribution contains 138 raw audiovisual source sessions: 69 with
the \texttt{P} prefix and 69 with the \texttt{PP} prefix. Video, audio, transcript and aggregate-
outcome records join one-to-one under the raw source-session identifier. Filename numbers are not
treated as intrinsic participant identity. The approved canonical-to-source mapping defines 69
participant pairs. Canonical \texttt{P32} loads source \texttt{P32}, canonical \texttt{PP32}
loads source \texttt{PP33}, canonical \texttt{P33} loads source \texttt{P33}, and canonical
\texttt{PP33} loads source \texttt{PP32}. The completed corpus-wide audit confirms that all other
rows are identity mappings. Raw identifiers and files remain
unchanged.

This correction was established before participant grouping, feature extraction, proxy
construction or paired analysis. The mapping preserves session phase and uniquely matches all
138 source recordings to the 69 participant pairs. The face audit covers 138 sessions and the
identity audit covers all 69 pairs. A valid primary-analysis row requires a
readable video, readable audio, non-empty supplied transcript, one aggregate rating row, its raw
source-session identifier and its canonical participant/session keys.

The primary target is the expectation--maximisation aggregate value for
\emph{RecommendHiring}:
\begin{equation}
    y_{is}=\frac{\text{AGGR RecommendHiring}_{is}-1}{6}.
\end{equation}
The documented 1--7 rating is therefore mapped to $[0,1]$. It is called
\emph{crowd-perceived hireability} throughout. Individual-rater rows and the other
rating dimensions may describe measurement reliability or support explicitly secondary analyses;
they do not replace the primary target after results are known.

Measurement quality is audited before scorer interpretation. Ordinal Krippendorff alpha is
computed over individual \emph{RecommendHiring} ratings with sessions as units and workers as
raters. The EM aggregate is compared with the unweighted per-session mean using Pearson and
Spearman association plus raw-scale absolute error. These checks describe target reliability; they
do not tune the target or model.

The supplied inventory establishes 138 matched audio/video records, 138 aggregate ratings and
138 non-empty transcripts. Before modelling, the data checks confirmed these joins, the 69
participant pairs and modality completeness. Any unresolved source discrepancy would exclude
the affected analysis row rather than be filled by assumption.

\subsection{Pre/post interpretation}

The pre/post indicator is not used to suggest randomisation. Primary models predict sessions from
both phases. Supplementary tables report paired participant differences in observed ratings and
locked OOF predictions. These estimates are descriptive associations. They cannot identify a
causal effect of counselling because the corpus has no randomised untreated comparison group and
time is confounded with intervention exposure.

\subsection{Permitted legacy features}

The supplied facial and SHORE tables have known missingness or truncation in some sessions, while
the prosody files require schema correction. They are excluded from primary model inputs. Their
permitted uses are (i) verifying duration and coverage, (ii) documenting quality failures, and
(iii) predeclared secondary explanations that are clearly separated from the fresh-encoder
results. In particular, they cannot be used as a fallback when a primary modality fails.

\section{Modality representation boundaries}

\begin{sloppypar}
Visual and text inputs are candidate-only. Audio is a bounded sample of mixed-speaker interview
context containing participant and interviewer. A single outcome-blind, candidate-presence-based
active interval governs visual and audio sampling.
\end{sloppypar}

\begin{description}
    \item[Visual input.] The full video recording is scanned at one
    hertz with the pinned SCRFD-10GF detector~\cite{Guo2022SCRFD}. At each second, the operational
    candidate is the unique most-central detection: its centre must lie within 0.35 frame diagonals
    of the frame centre and, when another face is present, must lead the runner-up by at least 0.05
    frame diagonals. The first valid second must occur within 30 seconds. Starting there, the scan
    closes after a run of 30 invalid seconds once at least 60 seconds have elapsed, or at the end of
    the stream; the half-open active interval ends one second after its last valid detection. At
    least 32 candidate-valid seconds are required. This is a deterministic appearance-and-position
    heuristic with sample-level quality assurance, not face recognition or identity tracking.

    All candidate-valid integer-second indices inside that interval are retained when there are at
    most 128. Otherwise, if their ordered list has length $N_v$, retain the values at ranks
    $\lfloor j(N_v-1)/127\rfloor$, $j=0,\ldots,127$. Only those validated crops enter DINOv2.
    Interviewer or background faces and full frames are never substituted. The selection record
    preserves the interval, accepted crops and reasons for any failure.

    \item[Interview audio.] The corpus channels are
    near-duplicates and are not speaker-isolated. The deterministic mono 16-kHz waveform is sliced
    to the exact face-derived half-open interval before WavLM processing, so it still contains both
    participant and interviewer. The active waveform is divided into non-overlapping 20-second
    windows, including a non-empty final partial window. All windows are retained up to 16;
    otherwise 16 are selected by the endpoint-rank rule over the ordered active-interval windows.
    The model therefore sees at most 320 seconds of standardised active-interview context. Full and
    active durations, interval samples, window indices and source-sample boundaries are stored,
    with no cross-window padding leakage. Exact FFprobe timestamp integers bind each supplied WAV
    to the embedded audio stream by rational cross-multiplication, and the decoded 16-kHz sample
    count must agree exactly. Only when the active interval closes at the final video grid cell may
    physical audio end slightly earlier: the shortfall must be strictly below one verified source
    audio timestamp tick (3/125 seconds, or 384 decoded samples) and is recorded without padding.
    Every other shortage is a structured failure. The modality is called \emph{interview audio}, never
    candidate voice. Verified candidate-speech segmentation, if later available, is an optional
    sensitivity analysis and cannot replace or tune the core result.

    \item[Text.] Exact audited speaker aliases define candidate and interviewer roles. A colonless
    pipe fragment inherits candidate status only from the most recent explicit candidate marker;
    fragments after interviewer/unclassified markers or before any role are excluded. The observed
    \texttt{r:} marker is counted as unclassified, and any new colon-labelled marker fails rather
    than being guessed. The current inventory retains 77 continuation fragments (18,214 source
    characters) and excludes 3 (522 characters).
\end{description}

A boundary audit samples sessions across both phases and records false-candidate inclusion,
candidate omission and unresolved identity. Boundary rules are fixed without outcomes or audit
attributes. Sessions that cannot satisfy a modality boundary receive a structured failure code;
primary results use the frozen common complete-pair participant manifest, with failure counts and reasons reported
without post-hoc imputation.

Before any model is fitted, an outcome-blind boundary pass establishes a common eligibility
roster. A participant is retained only when both sessions have
valid visual, interview-audio and candidate-text summaries. No failed modality is imputed and no
unpaired session is retained. The original 69-participant fold and random-proxy assignments are
subselected without reranking. Let $E$ denote the eligible complete-pair count; the report gives
$E$, retained counts per original fold, every exclusion and its structured reason.

\subsection{Frozen encoders}

The visual representation is the final-layer CLS vector from DINOv2-large~\cite{Oquab2024}.
The audio representation is the final hidden state of WavLM-Base+ at
16 kHz~\cite{Chen2022WavLM}. Each selected bounded session-audio window is encoded independently; valid frame
centres are pooled into consecutive zero-based one-second bins, retaining a non-empty final
partial second, and each bin concatenates mean and maximum values. Window summaries are then
aggregated deterministically without weighting padded samples. The text representation uses
MPNet pretraining~\cite{Song2020MPNet} within the Sentence Transformers
framework~\cite{Reimers2019}, specifically the
all-mpnet-base-v2 checkpoint~\cite{SentenceTransformersAllMPNet2025}.
Tokenizer IDs are divided into non-overlapping chunks of at most 382 content
tokens plus model special tokens, followed by masked mean pooling and L2 normalisation.

Encoder versions, preprocessing settings and output dimensions were fixed for all configurations.
Each source recording was checked against its approved identity and content before extraction;
completed representations were accepted on resumption only after the same checks passed. The
encoders remained frozen throughout model training.

\section{Session-level perceived-attribute annotations}

\subsection{Supplied source and construct}

The restricted annotation source contains 126 session-level rows supplied as one observer's
perceptions: all 69 post sessions and 57 of the 69 pre sessions. The source rows were preserved
and matched to the verified participant/session mapping without altering the observer's codes. The
annotations were supplied only after the Stage-A predictions and proxy assignments had been
locked, so they could not affect ordinary fitting or model selection.

The accepted visual constructs are:
\par\noindent\begin{minipage}{\linewidth}
\begin{enumerate}
    \item perceived gender presentation, preserving the observed source codes \texttt{F} and
    \texttt{M} as feminine and masculine presentation;
    \item perceived racialised appearance, preserving the source categories and treating
    \texttt{CD} (cannot determine) as unknown; and
    \item perceived skin tone under the recorded imagery and lighting on the integer ordinal
    scale from 1 (lightest) to 5 (darkest).
\end{enumerate}
\end{minipage}\par
These labels describe one observer's perception in a particular recording. They are never
described as true gender, race or ethnicity, and the source column headed ``Ethnicity'' is renamed
only as perceived racialised appearance in analysis. Self-reported attributes, if later
authorised, remain a separate construct and are never combined with these visual judgements. The
source age-band values are retained only for provenance and reconciliation; age is excluded from
Stages B and C.

The source does not provide a reproducible rater identifier, field-level confidence, collection
instructions or a blinding record. These facts remain unknown and are reported as limitations.
No independent second judgement exists, so the study makes no inter-rater agreement or consensus
claim and does not manufacture either from the pre/post rows.

\subsection{Phase-aligned primary and sensitivity analyses}

The primary Stage-C audit joins each of the 69 locked post-session OOF predictions only to its
matching post-session annotation. The 57 available pre annotations define a separate,
reduced-coverage sensitivity: each is joined only to its matching pre-session prediction. A post
annotation is never attached to a pre prediction, and pre/post values are neither averaged nor
adjudicated into a participant consensus. Post is primary because it supplies a common complete
roster, not because a post judgement is assumed to be more accurate or stable.

For the 57 participants with both rows, an additional stability diagnostic reports exact
pre/post matches for the nominal fields and the absolute pre/post difference distribution for the
ordinal skin-tone field. This is an intra-source, across-session sensitivity to presentation,
recording and phase; it is not inter-rater reliability, and it cannot establish that either value
is a stable identity attribute. The incomplete pre roster is reported explicitly and never
substitutes for the complete post primary audit.

\subsection{Sparse-group rule}

A category is reportable only if it has at least 10 unique participants in the relevant
phase-aligned roster. No category is merged with another after outcomes or predictions are
viewed. An axis with fewer than two reportable groups receives counts and coverage only; no
disparity estimate is produced. Intersectional analysis uses the same 10-participant minimum and
is secondary. Unknown responses are reported as missing audit information, never treated as a
demographic group in the headline comparison. Perceived skin tone remains ordinal 1--5 and is
never thresholded or binarised to increase support.

\subsection{One-way firewall}

Audit annotations may be collected early, but are stored separately, unreadable by ordinary
loaders and not inspected by model developers until the ordinary scorer artefacts are locked. The
analysis proceeds through exactly three irreversible stages:
\begin{enumerate}
    \item[Stage A.] Evaluate nine deterministic reference models and seven learned
    configurations (105 fits). Preserve their held-out predictions, participant folds and proxy
    assignments before opening the audit annotations. No perceived-attribute label is used.
    \item[Stage B.] Use the post-session \texttt{F}/\texttt{M} field for one exploratory,
    segregated perceived-gender-presentation oracle. Although the source legend also
    permits \texttt{A} (androgynous or ambiguous) and \texttt{CD} (cannot determine), every
    supplied row is coded \texttt{F} or \texttt{M}; no observed
    category is dropped or merged. This observed-data binary axis was chosen after annotation
    receipt, not before rating, and is frozen before any Stage-B fit or Stage-C audit. Within each
    rotation, the oracle may use the one post-derived label for each training participant, but no
    held-out-participant annotation. It runs only if both training groups contain at least 10
    participants; otherwise it records a zero-fit omission.
    \item[Stage C.] In a read-only evaluator, join post predictions to the 69 post labels for the
    primary audit and pre predictions to the 57 available pre labels for the sensitivity. Then
    permit no model rerun, proxy choice, hyperparameter change or selective omission in response
    to the audit; any later analysis is labelled post hoc.
\end{enumerate}

\section{Fixed cross-fitting design}

Participants, not sessions, define folds. A fixed, outcome-independent ordering of the corrected
participant identifiers assigns each person to one of five folds, with sizes 14, 14, 14, 14 and
13. Both pre- and post-intervention sessions follow that assignment, irrespective of the raw
recording number. The assignments were fixed before outcomes were used for training and before
perceived-attribute annotations were examined.

For outer fold $f$, fold $f$ is test, fold $(f+1)\bmod 5$ is validation, and the other three folds
are training data. For each configuration and seed $r\in\{1,11,21\}$, the model predicts only the
outer test fold. Training-only means, scales, proxy transforms and cluster centres are estimated
on the three training folds; stopping uses the fixed validation fold. The five test parts form one
OOF prediction for
every eligible session; no in-fold prediction enters a primary metric. With five folds and three
seeds, each learned configuration has 15 fits. Each configuration--seed pair has exactly $2E$
OOF session predictions, or $6E$ prediction rows across three seeds; these totals are 138 and 414
only if $E=69$. The three seed-specific OOF vectors are retained, not collapsed before
diagnostics.

There are exactly seven learned Stage-A configurations, hence $7\times5\times3=105$ learned
fits. Exactly nine deterministic cross-fitted diagnostic floors are outside this matrix and add no
seeded learned fits: \texttt{mean}, \texttt{phase\_mean}, \texttt{ridge\_video},
\texttt{ridge\_audio}, \texttt{ridge\_text}, \texttt{ridge\_video\_audio},
\texttt{ridge\_video\_text}, \texttt{ridge\_audio\_text} and \texttt{ridge}.

The original experiment contains no broad hyperparameter search on its 69 participants.
Architecture and optimisation settings are frozen through synthetic tests and outcome-blind
engineering checks. Ordinary ERM shuffles complete participants deterministically and uses full
16-participant batches followed by the final complete-participant remainder; it never drops or
pads that remainder, and both sessions of each participant remain together. On the complete
69-participant fold manifest the training schedule is therefore 16/16/9 or 16/16/10 participants,
depending on the rotation. Proxy-balanced configurations retain their separately fixed
16-participant cyclic policy. Every learned configuration uses AdamW with learning rate
$10^{-3}$ and weight decay $10^{-2}$, at most 200
epochs, gradient-norm clipping at 1.0, and patience 25 on validation SRCC. The maximum-validation-
SRCC checkpoint is retained, with the earliest epoch resolving an exact tie. Linear weights use
Xavier-uniform initialisation, linear biases are zero, and LayerNorm scale/bias begin at one/zero.
The W1 coefficient is fixed at $\lambda_W=0.02$ from the first parameter update and is never
swept or changed using audit information.

\section{Shared low-capacity predictor}

For each session, valid DINOv2 face vectors, bounded-window WavLM summaries and MPNet chunk
vectors are averaged within modality. A separate standard scaler is fit on training participants
only. Each modality then passes through \texttt{Linear}$(d_m,8)$, \texttt{LayerNorm}$(8)$ and a
GELU activation. The three eight-dimensional outputs are concatenated into 24 dimensions and
passed through \texttt{Linear}$(24,1)$ and a sigmoid. The 26,721 trainable parameters are
deliberately modest for 69 independent participants; encoder weights remain
frozen. This late-fusion network is used unchanged for ERM, sampling controls, random-proxy W1,
visual-proxy W1, audio-proxy W1, ARL and the oracle; only the declared sampler or loss term changes.
The \texttt{mean} floor predicts the rotation's arithmetic training-target mean, while
\texttt{phase\_mean} predicts separate arithmetic means for canonical pre and post training rows.
Seven participant-safe ridge floors use video, audio, text, each pair, and all three modalities,
respectively. Within each rotation, each ridge concatenates only its named 1024-dimensional
DINOv2, 1536-dimensional WavLM and/or 768-dimensional MPNet session summaries. A StandardScaler is
fitted on training session rows, followed by scikit-learn Ridge on the normalised target with
\texttt{alpha=1}, an intercept, \texttt{solver=lsqr}, tolerance $10^{-6}$ and 10,000 maximum
iterations. The unchanged scaler and model are applied to held-out rows; there is no tuning or
validation refit.

The base loss for session $j$ is mean squared error,
\begin{equation}
    \mathcal{L}_{\mathrm{ERM}}=\frac{1}{m}\sum_{j=1}^{m}(y_j-\hat y_j)^2.
\end{equation}
These optimisation, regularisation, batch-construction and stopping settings are recorded in every
run manifest. The sigmoid head already constrains predictions to $[0,1]$; no additional clipping
or post-processing is selected from validation or audit results.

\section{Proxy construction}

\begin{sloppypar}
Visual proxies are constructed from participant-level DINOv2 candidate-face aggregates across both
available sessions within the applicable fold; audio proxies are constructed analogously from
participant-level active-interval interview-audio WavLM aggregates across both sessions. For each source, a
training-only standard scaler followed by scikit-learn k-means is fit in ascending numeric
participant ID order, using two clusters, k-means++ initialisation, 20 initialisations, 300 maximum
iterations, tolerance $10^{-4}$, Lloyd's algorithm and random state 0. Both sessions contribute to the participant representation only within the same training
fold, and each participant receives one proxy assignment. Held-out assignments use the frozen
nearest-centre rule. Cluster identifiers are arbitrary and do not denote a perceived attribute.
\end{sloppypar}

Proxy-balanced sampling draws eight participants from each available proxy group into a
16-participant batch and includes both sessions, yielding 32 rows. Both groups are required. A
deterministic, seed-controlled cyclic reshuffle is used within each
eight-participant half; duplication within a half occurs only when its training group has fewer
than eight participants. Participant-level unique and duplication rates are logged. The
sampling-only controls use this policy with the ERM loss and do not pass proxy groups to the
objective. If a rotation contains $N_{\mathrm{train}}$ training participants, every proxy-
balanced configuration performs exactly $\lceil N_{\mathrm{train}}/16\rceil$ optimizer updates
per completed epoch. A single seeded cyclic sampler state advances continuously across epoch
boundaries, and every draw and repeat is persisted. Validation and test use their complete natural
role distributions rather than balanced resampling. The
W1 configurations add
\begin{equation}
    \mathcal{L}=\mathcal{L}_{\mathrm{ERM}}+\lambda_W
    W_1\!\left(\{\hat y_j:g_j=0\},\{\hat y_j:g_j=1\}\right),
\end{equation}
using the same corresponding balanced sampler. The coefficient $\lambda_W=0.02$ is applied exactly
once from the first update and is fixed before the annotation audit. The proxy assignments,
training schedule, fitted scalers and model checkpoints were preserved for the subsequent
comparison.

The random control uses one fixed, outcome-independent participant ordering: the first 34 form
group 0 and the remaining 35 group 1. This assignment is
identical across both sessions and all fold rotations, and is independent of phase, outcome,
representations and audit annotations.

Adversarially reweighted learning (ARL) adapts the example-reweighting idea of
Lahoti et al.~\cite{Lahoti2020} to continuous regression. Putra et al.\ applied a related
approach to interview videos~\cite{Putra2024}; the present model is not an exact reproduction
of their system. Let
$h_j\in\mathbb{R}^{24}$ be the common scorer's pre-head representation and let $y_j$ be the
normalised continuous target. A training-only affine adversary with 26 parameters receives
$[\operatorname{stopgrad}(h_j),y_j]$ and produces $q_j=\operatorname{sigmoid}(a_\phi(\cdot))$.
For the current minibatch of $m$ session rows, its learner weight is
\begin{equation}
    w_j=1+\frac{q_j}{\max\!\left(m^{-1}\sum_{k=1}^{m}q_k,10^{-4}\right)},
    \qquad
    \mathcal{L}_{\mathrm{ARL}}=\frac{1}{m}\sum_{j=1}^{m}
    \operatorname{stopgrad}(w_j)(y_j-\hat y_j)^2.
\end{equation}
The loss is a batch mean and is not renormalised by the weight sum; the full 32-row batch has
upper bound $w_j\leq33$. The first six learner minibatch updates use ordinary MSE and do not call
or step the adversary. From update seven onward, the learner is updated first; current detached
representations and squared errors are then recomputed, and the adversary alone minimises the
negative weighted mean error. Learner and adversary use separate AdamW optimisers with the common
learning rate and weight decay and separate norm-1 gradient clipping. ARL retains the ordinary
participant-paired sampler, including its complete final remainder, and the unweighted validation-
SRCC checkpoint rule. Each checkpoint stores both models, both optimiser states, the train-only
scalers, training phase and update counts. Canonical minibatch and epoch records persist the
objectives; adversary-logit, sigmoid-score, guarded-denominator and weight summaries; effective
sample size; weight--error Pearson correlation; mean absolute within-participant-pair weight
difference; pre-clipping gradient norms; and unweighted validation SRCC and MSE. Statistics
undefined during learner-only pretraining or for a constant vector are stored as JSON null rather
than NaN. Putra et al.'s combined MSE+BCE variant, not its MSE-only variant, was its strongest
reported First Impressions system; consequently this architecture-controlled MSE comparator is
not described as reproducing that result.

The corpus-annotated oracle replaces the proxy group in the W1 loss with the supplied
post-session perceived-gender-presentation field, using one post-derived label for each training
participant. All observed values are \texttt{F} or \texttt{M}, mapped to feminine and masculine
presentation. This is a binary coding of the observed source rows, not a claim that gender itself
or the source legend is binary. The axis was selected after the annotations were received and is
therefore an exploratory comparison; it was not preregistered or designated before rating. It is
nevertheless frozen before any Stage-B fit or Stage-C audit. Self-reported and perceived
provenance are never pooled participant by participant.

The oracle is not a deployable demographic-free method; it is a segregated diagnostic of whether
weak proxy--audit alignment can explain a result. If either perceived-presentation group has fewer
than 10 training participants in any rotation, the oracle is omitted by rule and the omission is
reported. No alternative axis is substituted and no category is collapsed merely to make the
oracle run. Perceived skin tone is restricted to Stage C and remains ordinal 1--5.

\section{Comparators and contrasts}

\begin{table}[htbp]
\centering
\caption{Predeclared configurations and their inferential roles.}
\label{tab:configurations}
\begin{adjustbox}{max width=\textwidth}
\begin{tabular}{lll}
\toprule
Configuration & Training signal & Role \\
\midrule
Nine fixed floors & mean/phase mean and seven ridge modality subsets & deterministic diagnostics \\
ERM & MSE, ordinary participant sampling & common-backbone reference \\
Visual sampling & visual-proxy balanced sampling, MSE & sampling control \\
Audio sampling & audio-proxy balanced sampling, MSE & sampling control \\
Random W1 & random groups, balanced sampling, W1 & generic penalty control \\
Visual W1 & visual proxy, balanced sampling, W1 & primary intervention \\
Audio W1 & interview-audio proxy, balanced sampling, W1 & primary intervention \\
ARL & adversarial example weights & generic demographic-free comparator \\
Oracle W1 & training-fold audit group, W1 & segregated explanatory comparator \\
\bottomrule
\end{tabular}
\end{adjustbox}
\end{table}

The primary total-effect contrasts are Visual W1 minus ERM and Audio W1 minus ERM. Incremental W1
contrasts are Visual W1 minus Visual sampling and Audio W1 minus Audio sampling. Random W1 minus ERM
tests a generic partition penalty, ARL minus ERM tests a generic demographic-free method, and
Oracle W1 minus ERM describes the label-supervised reference where estimable. A mitigation
interpretation requires each eligible contrast to be assessed for both utility and audited
disparity; conclusions cannot be based on one axis alone.

\section{Outcomes and uncertainty}

\subsection{Utility}

Primary utility is Spearman rank correlation between all eligible session-level OOF predictions
and aggregate \emph{RecommendHiring}. Mean absolute error is co-primary descriptive support, and
root mean squared error and Pearson correlation are secondary. Metrics are computed separately
for seeds 1, 11 and 21;
the three OOF vectors are not averaged into an ensemble. Paired model contrasts are formed within
seed and their arithmetic mean is the declared summary effect. A constant prediction within one
test fold has undefined within-fold rank correlation. The cross-fitted mean floor can vary across
outer training folds, so its pooled held-out correlations are reported alongside its errors.

\subsection{Audited disparity}

For each reportable phase-aligned audit axis, the disparity statistic is the maximum pairwise
empirical W1 distance between the score distributions of the single-observer perceived groups:
\begin{equation}
    D_{\mathrm{audit}}=\max_{a<b}W_1(\hat Y_a,\hat Y_b).
\end{equation}
Group means, medians, quantiles and counts accompany this maximum so that it is not interpreted
without scale or composition. Direction is reported for two-group axes; multi-group results name
the pair attaining the maximum. Observed target disparity is reported beside prediction disparity
to distinguish reproduction of the rating distribution from model-added disparity.

\subsection{Participant-cluster intervals}

Uncertainty respects the participant as the independent unit. For the primary Stage-C audit, a
fixed 10,000-replicate bootstrap with seed 2026 samples the 69 post-phase participants with
replacement and retains each sampled participant's matching post prediction and annotation. The
pre sensitivity independently resamples its 57 available participants and retains each matching
pre prediction and annotation. Each replicate recomputes group disparity, every within-seed paired
configuration difference and their arithmetic mean on locked OOF predictions. Percentile 95\%
intervals and the fraction of valid replicates are reported. A replicate in which a required group
is absent is invalid for that attribute only and is never imputed. Inference emphasises effect
sizes, interval width and consistency across seeds rather than dichotomous significance tests.

\subsection{Descriptive paired-change accuracy}

For participant $i$, define the observed and predicted changes from the two OOF sessions as
\begin{equation}
    \Delta y_i=y_{i,\mathrm{post}}-y_{i,\mathrm{pre}},\qquad
    \Delta\hat y_i=\hat y_{i,\mathrm{post}}-\hat y_{i,\mathrm{pre}}.
\end{equation}
The secondary paired-change package reports Spearman correlation between these $E$ deltas, delta
MAE $E^{-1}\sum_i|\Delta y_i-\Delta\hat y_i|$, and signed delta bias
$E^{-1}\sum_i(\Delta\hat y_i-\Delta y_i)$. These are computed from pooled OOF predictions, never
in-fold fits, and receive participant-bootstrap intervals. They measure descriptive change
tracking only; they do not estimate a coaching effect.

\section{Explanatory and sensitivity analyses}

The fixed failure-analysis package contains:
\begin{enumerate}
    \item proxy size, stability across outer rotations, byte identity across seeds, centre
    separation and visual--audio agreement;
    \item post-firewall proxy--audit association using contingency tables, adjusted Rand index and
    normalised mutual information, interpreted only as association;
    \item inference-only modality masks on each locked ERM model: set one modality's
    training-standardised summary to zero (the training mean), with no refitting or retuning;
    \item signed-residual and absolute-error summaries by reportable audit group;
    \item phase-stratified summaries and paired pre/post delta correlation, MAE and signed bias;
    \item complete-pair eligibility coverage and structured modality-boundary failure summaries; and
    \item secondary comparisons with legacy-feature quality indicators, not primary replacements.
\end{enumerate}
These analyses explain mechanisms and boundary failures. They cannot select the preferred proxy,
change sparse-group thresholds or rescue an unfavourable primary result.

\section{Exploratory linear follow-up on MIT}
\label{sec:linear-followup-method}

The original scorer's weak predictive performance limits what its penalty comparison can
explain. The follow-up therefore uses the same accepted three-modality summaries and
participant folds with an identity-output linear scorer. This gives 3,328 input coefficients
and an intercept. A training-row standard scaler is fitted within each outer rotation.
The natural task-only objective is the alpha-1 ridge objective, with the intercept left
unpenalised. The model is initialised at its exact task-only solution rather than at a new
zero or random predictor. Predictions are not clipped to the target range.

For training rows $j=1,\ldots,n$, the common objective is
\begin{equation}
 J_q(w,b)=\sum_{j=1}^{n}q_j\bigl(y_j-x_j^{\mathsf T}w-b\bigr)^2
       +\frac{\|w\|_2^2}{n}
       +\lambda W_1(\hat Y_{g=0},\hat Y_{g=1}).
\end{equation}
Natural exposure uses $q_j=1/n$. Deterministic balanced exposure uses
$q_j=1/(2n_{g_j})$, where $n_{g_j}$ is the number of session rows in that proxy group.
The latter is the expected group-balanced task weighting; it is not an exact replay of
the original cyclic minibatch sampler. Each participant's two sessions share one group.
Visual and audio proxies retain the original training-only participant-mean clustering
policy. No outcome or perceived attribute selects the proxy partition.

The full training-fold W1 uses the integral of the empirical step quantile functions,
including unequal group sizes. Linearly interpolating two sorted samples onto a shared grid
would generally produce a different statistic and is not used. Optimisation takes place in
the training design's singular-vector row space, preserving the minimum-L2 linear model
class. L-BFGS-B receives the analytic score subgradient, with derivatives averaged within
exact score ties. The fixed cap is 1,000 iterations. Since projected-score W1 is not globally
convex in the coefficients, termination is reported without a global-optimum claim.

The comparison contains one natural task-only reference, visual and audio balanced task-only
references, and three strengths ($0.002$, $0.02$, $0.2$) for each proxy under natural and
balanced exposure. These 15 deterministic configurations run in all five outer folds,
giving 75 fits. Repeating identical full-fold fits under nominal seeds would not provide
additional independent evidence. Each penalty is compared with its own zero-penalty exposure
control, and every strength is reported rather than selecting a winner from test or audit
results. The four contrasts at $\lambda=0.02$ are the primary descriptive comparisons.

All predictions are fixed before joining the accepted phase-specific perceived annotations.
The 69-post/57-pre audit and sparse-category rules remain unchanged. The manipulation check
computes proxy W1 separately within each outer test fold, with that fold's fitted group
assignments; neutral cluster IDs are not stitched into a global group identity. Utility and
audited-disparity contrasts use 10,000 paired participant bootstrap draws, conditional on
the fitted models. Both sessions stay together for pooled utility. The post and pre audits
resample only the matching phase. All numerical terminations remain visible.

\subsection{Outcome-only ridge selection analysis}

A separate predictive comparison examines whether regularisation selection improves the
benchmark without consulting perceived attributes. Within each outer rotation, its three
training folds supply three inner participant-grouped rotations. All scalers are fitted
only on the corresponding inner training people. Ridge alpha is chosen from
$\{0.1,1,10,100,1000\}$ by pooled inner absolute error, with smaller alpha for an exact tie.
The broader comparison also includes a ridge-objective-matched linear candidate, a small
late-fusion neural candidate and the fixed reference predictors. No test outcome selects
the candidate or alpha, and the reserved outer validation fold is not added to training.

The available selected-ridge summary is an unweighted mean of the five outer-fold metrics.
It is reported with that label. In particular, a mean of fold-specific Spearman correlations
is not the pooled out-of-fold correlation used for the original and compact comparisons.
The summary does not supply a participant-bootstrap interval or a paired comparison with
those pooled estimates, so neither is inferred from the point values.

\section{AVI-Personality prediction and error audit}
\label{sec:avi-method}

AVI-Personality supplies six recorded answers per participant and participant-level
recruiter ratings~\cite{Zhang2026AVI}. The validated feature release covers all 646 people.
The official splits contain 452 training, 64 validation and 130 test participants. Two
training participants lack the primary interview-performance outcome, leaving 450 people
for supervised fitting; they remain represented in source missingness and feature coverage.
No person or question clip crosses an official split.

The visual input differs from the MIT candidate-face boundary. DINOv2 encodes full frames
sampled at one hertz, capped at 128 per question clip. Their mean forms a 1,024-dimensional
clip vector. Embedded audio is decoded to mono 16 kHz and divided into non-overlapping
20-second windows. WavLM representations are mean--maximum pooled into one-second bins,
then averaged into a 1,536-dimensional clip vector. A final residual shorter than the
encoder's 400-sample receptive field is omitted rather than padded into invented content.
Each modality's six question vectors receive equal weight in the participant mean.
The modelling release requires all six valid clips. No transcript representation is supplied,
and no zero text vector is created to imitate a third modality.

The predictive baseline concatenates the visual and audio participant vectors and standardises
each feature using training participants only. A ridge scorer selects alpha from
$\{0.1,1,10,100,1000\}$ by validation MAE, preferring smaller alpha on an exact tie.
It predicts the original 1--5 recruiter rating and clips predictions to that range.
The training-target mean is a separate reference. A preliminary division of modality blocks
by their dimensional scales is cancelled by subsequent featurewise standardisation;
the executed model is therefore described as standardised concatenated ridge, without
claiming a distinct modality-balancing effect.

After predictions are fixed, a separate audit joins source self-reported gender. The test
contains 67 male and 62 female respondents, one prefer-not-to-say response and no non-binary
or third-gender respondent. The two supported groups receive descriptive MAE estimates;
the nonresponse is counted without forming a one-person disparity estimate. The signed
male-minus-female MAE gap and its absolute magnitude describe prediction error, not score
parity or the effect of a fairness intervention. A percentile interval obtained by
bootstrapping the absolute gap is not used as a test of a nonzero signed population
difference. Self-report is not substituted for the MIT observer's perceived presentation.

This independent baseline estimates prediction and group-conditioned error in another
corpus. Considered alone it contains no W1-versus-control comparison, and the matched
comparison described next was specified only after these baseline results had been
inspected. Outcome scales, visual boundaries and rating sources
differ, so raw scores or errors are not pooled across the two corpora. Any method
comparison designed after inspecting these test results is exploratory on independent
people, rather than an untouched confirmatory holdout.

\subsection{Matched intervention and control comparison}
\label{sec:avi-compact-method}

A subsequent comparison applies the proxy-regularisation design of the main study to these
independent participants. Because the predictive baseline and its group-error audit above
had already been inspected, this comparison is exploratory on independent people rather
than an untouched confirmatory holdout. Its purpose is to establish whether the proxy
penalty behaves on a second corpus as it did on the first, not to confirm an effect.

Fifteen fixed settings are fitted. Three task-only controls carry no penalty: a
natural-exposure control, and one balanced-exposure control for each proxy source. Twelve
penalised arms cross two proxy sources, visual and audio, with two exposure policies,
natural and balanced, and the three penalty strengths $\lambda\in\{0.002,0.02,0.2\}$ used
in the main study. Each penalised arm is paired with the matched control sharing its
exposure policy, so a penalty effect is never read against an unmatched reference. The
complete grid is reported; no strength is selected using test or audit results.

Proxies are obtained without protected labels, by the same unsupervised partition of
frozen representations used in the main study. The visual partition assigns 368 training
and 103 test participants to its larger group against 82 and 27 respectively; the audio
partition assigns 143 training and 35 test participants to its smaller group against 307
and 95. Ridge regularisation is selected once on validation error and held fixed across
all arms, so arms differ only in exposure policy, proxy source and penalty strength.

All held-out predictions for all fifteen arms are locked before the separate source gender
file is opened, so audit attributes cannot influence fitting or model selection. The locked
test set contains 130 participants and 1{,}950 prediction rows. Paired uncertainty uses
10{,}000 participant-level bootstrap draws, with the same draws applied to every arm so
that matched contrasts share resampling noise. The resulting intervals are pointwise and
conditional on these fitted models; they describe sampling variation in the evaluated
predictions rather than uncertainty over refitting, and they are not corrected for
inspecting twelve penalised arms across several outcome measures.

Three quantities are reported for each matched contrast: the change in original-scale
mean absolute error, the change in Spearman rank correlation, and the change in the
Wasserstein-1 distance between the predicted-score distributions of the two supported
source gender groups. The last is reported as a disparity gain, where a positive value
indicates a smaller score gap than the matched control. Held-out proxy-group W1 is
reported separately from audited gender disparity, because a penalty acting on an
unsupervised partition is not thereby acting on a protected attribute.

\section{Reproducibility and data protection}

The study preserves its source inventory, participant folds, fixed encoder and model settings,
training seeds and held-out predictions for verification. The raw recordings, identifiable face
crops and individual annotations are not redistributed. Published results are aggregate, and
the analysis of perceived attributes remained separate from ordinary model development.
